% Pursuing Stacks (cotes 134-2 to 134-8): archive cote 134-4, site page 39 (PDF page 40), typescript page 127.
% typescript. Draft compared: ps3.tex:1948-2013.
% AI-transcribed by Claude Opus 5.5 (claude-opus-5-5, adaptive thinking, effort high), Message Batches API, batch ps2 (msgbatch_01NVwr1Yuo8qJf5rykeGBXPK), 2026-11-12; not yet checked by eye.
\scanpage{134-4/39}
\mspage{127}
where now
\[ i\colon M \hand{\to} W^{\hand{-1}}M = M_W \]
is the canonical functor defined in terms of W;  we thus get a canonical \unsure{f}
functor of categories ``under M''
\[ (\over{2}{3}) \qquad \hand{\beta=\beta_W}\colon M_R \hand{\to} M_W \quad , \]
as a matter of fact,  R is the largest element in Q (for the inclusion re-%
lation of subsets of $F\ell(M)\mathrm{x}F\ell(M)$) for which a functor (\over{2}{3}) exists (then ne-%
cessarily unique,  as before).  We may say that W is a homotopism structu-%
re on M,  i.e.  $W\hand{\in}L_0$  (where $L_0$ is the subset of L of all homotopism
structures on M) iff the functor (\over{2}{3}) is an isomorphism,  or equivalently
(as it is clearly bijective on objects,  \over{\unsure{sur}}{inj}ective on arrows)  iff it is
\uline{surjective on arrows} .

We may describe $Q_0$ and $L_0$ in a purely set theoretic way,  in terms
of the system (r,s) of maps in (1),  by the formula (which is just a
translation of the definitionsof $Q_0$ and $L_0$)
\[ Q_0 = \hand{\Big\{} q \hand{\in} Q \hand{\big|} sr(q) = q \hand{\Big\}} \]
\[ L_0 = \hand{\Big\{} \ell \hand{\in} L \hand{\big|} rs(\ell) = \ell \hand{\Big\}} \quad , \]

and we can describe formally the pair of subsets $(Q_0,L_0)$  of Q,L  as the
largest pair of subsets, such that r and s induce between $Q_0$ and $L_0$ bijec-%
tions inverse of each other.  In \unsure{this} general set-theoretic set up,  it is
by no means clear,  and false in general,  that r maps Q into $L_0$ or L\deleted{\unsure{$_0$}} int
% NOTE: the line runs off the right edge of the scan after "int"
$Q_{\hand{0}}$ (thus both $Q_0$ and $L_0$ may well be empty,  whereas Q and L are not).  Thu\unsure{s}
it is not clear at all that starting with an arbitrary $R\hand{\in}Q$,  the correspon-%
ding W = r(R)  is a homotopism structure,  and it is easily seen that this
is \uline{not} so in general,\deleted{\illegible{1}} contrarily to what I hastily stated in yesterday\hand{'}s
notes.  (Take M with just one object,  therefore defined by a monoïd G,
and $M_R$  corresponds to a quotient monoïd \inserted{$G'$}, we \uline{may} take \over{\illegible{1}}{$G'$}= unit monoïd,
thus W is G itself, and $M_W$ the enveloping group \inserted{$\overline{G}$} of G  -  the \xout{\illegible{1}} map
$\overline{G} \hand{\to} \over{\illegible{1}}{G'}$  need not be \over{\unsure{sur}}{in}jective ! )\hand{.}In the opposite direction,  starting with
an arbitrary W in L,  the corresponding R need not be a homotopy relation.
(If M is reduced to a point,  this amounts to saying that if we localize
a monoïd G with respect to a subset $W\hand{\subset}G$ ,  by making invertible the ele-%
ment in W,  the corresponding map $G \hand{\to} \overline{G}$ need not be surjective !)

What we may say,  though,  is that if we start with a pair
\[ (R,W) \hand{\in} Q\mathrm{x}L \]
such that the \uline{two} functors (2),(3)  exist,  i.e. such that
\[ W \hand{\leq} r(R) \quad\text{and}\quad R \hand{\leq} s(W) \quad , \]
then \xout{\illegible{1}} $(R,W)\hand{\in} Q_0\mathrm{x}L_0$  i.e. R is a homotopy relation and W is a homotopism
\margin{39}
